\blockhead{opEdge}{Block 10 --- Stage 5 operator
$\mathrm{ConvT}^{64 \to 1}_{k=4,\,s=2,\,p=1} \circ \tanh(\cdot)$}

The terminal stage differs from the previous four in two respects: (i)
it collapses the 64-channel feature map to a single output channel,
matching the grayscale nature of the target recurrence plot; (ii) it
replaces $\mathrm{BN} \circ \sigma_{\mathrm{ReLU}}$ with a pointwise
$\tanh(\cdot)$, which bounds the per-pixel output to the open interval
$(-1, 1)$. The $\tanh$ choice is consistent with the per-bucket
min--max normalisation applied to the training tensors, which scales
every training recurrence plot to $[-1, 1]$ via an invertible
$(\min, \max)$-tuple. Because the generator's output is bounded by the
same interval as the training data, the per-pixel histograms of
$G_\theta(z)$ and a training tensor live on the same support, which is
a precondition for the Wasserstein-1 distance estimator to be
well-defined under the W-GAN-GP duality. Note also that there is no
batch-normalisation at this stage: applying BN immediately before
$\tanh$ would impose a per-channel zero-mean constraint on the output,
biasing the network away from high-amplitude plots even when the
target distribution favours them.
